% !TeX root = ../practicat6t7t8.tex
% !TeX encoding = utf8

\chapter{Resumen}

El presente trabajo analiza las políticas de Formación y Desarrollo
(Tema~6), Evaluación y Gestión del Rendimiento (Tema~7) y Retribución
(Tema~8) aplicadas en CaixaBank, S.A., primer grupo bancario español
por volumen de activos, con más de 45.000 empleados.

En el ámbito de la formación, se examina el ciclo completo del programa
formativo de la entidad —identificación de necesidades, diseño,
implantación y evaluación—, así como el tiempo y coste asociados.
Se analizan igualmente las nuevas tendencias adoptadas: la plataforma
de \textit{e-learning} Virtaula, el modelo 70-20-10, la estructura de
universidades corporativas (escuelas especializadas internas),
el \textit{outdoor training} y la gamificación. Respecto a la gestión
de la carrera profesional, se describe el proceso de planificación
en sus cuatro fases —premisas, valoración, dirección y desarrollo—,
detallando las herramientas de rotación de puestos, \textit{coaching},
\textit{mentoring} e \textit{mentoring} inverso empleadas por la entidad,
junto con las tendencias de \textit{networking}, empleabilidad y marca
personal.

En el ámbito de la evaluación del rendimiento, se estudia la delimitación
de la evaluación (qué, cuándo y quién evalúa), los métodos objetivos,
subjetivos y mixtos utilizados, así como las estrategias de reforzamiento
positivo, negativo y la entrevista de evaluación. Se presta especial
atención a las tendencias de evaluación continua y por compromisos que
CaixaBank está implantando en sustitución de los sistemas anuales
tradicionales.

Finalmente, en materia retributiva, se analiza la estructura salarial de
cinco puestos representativos con referencia al Convenio Colectivo
sectorial de Banca, los tres tipos de equidad (interna, externa e
individual), el modelo híbrido de retribución por puesto y por
competencias, los principales incentivos financieros y elementos no
financieros, y las tendencias de retribución flexible y salario emocional.

La metodología combina el análisis de fuentes documentales públicas
—memoria anual, informe de sostenibilidad y Plan Estratégico
2025--2027 de CaixaBank— con una entrevista semiestructurada a una
directora de sucursal de la entidad.

\medskip

\textbf{Palabras clave:} formación corporativa, \textit{e-learning},
gestión de carrera profesional, \textit{coaching}, \textit{mentoring},
evaluación del rendimiento, retribución flexible, salario emocional,
CaixaBank, convenio colectivo de banca.

\cleardoublepage
\endinput
